
\chapter*{Nota editorial}
\addcontentsline{toc}{chapter}{Nota editorial}
Esta muestra aplica al nuevo título la maquetación de los volúmenes finales de la colección Studium: página A4, tipografía académica, márgenes, jerarquía de títulos, cabeceras, ecuaciones y cajas didácticas. La portada incorpora una fotografía del núcleo de un motor GEnx, con el crédito y la licencia indicados al pie de esta página.

El contenido siguiente permite comprobar el comportamiento de una apertura de capítulo, una página de desarrollo, un cuadro, una figura vectorial y un ejemplo con ecuaciones. Es una muestra parcial del futuro manual y no constituye la edición completa de la asignatura.

\caja{Estado de la edición}{\estado\ La revisión de maquetación no equivale a una auditoría científica o normativa del libro. El desarrollo del temario, la integración de Studium y las verificaciones del encargo general siguen pendientes.}

La titulación y la universidad identifican el contexto de estudio. Esta es una elaboración independiente y no se atribuye carácter oficial ni aprobación institucional.

\section*{Convenciones de lectura}
Las definiciones y los modelos se presentan en azul; los ejemplos resueltos, en ocre; las comprobaciones y síntesis, en violeta; las pautas de estudio, en verde. Cada ecuación, figura y cuadro tiene numeración de capítulo y puede citarse desde el texto.

Se emplean unidades SI. Los valores de los ejemplos son ficticios; permiten comprobar el procedimiento y la composición tipográfica. Los procedimientos reales de mantenimiento requieren documentación aplicable y controlada.

\par\vfill\creditosportada
\tableofcontents
\mainmatter
\chapter{Del material al comportamiento en servicio}
\label{ch:material}
\objetivos{Relacionar las condiciones de servicio con los mecanismos de daño y formular preguntas útiles para la selección de materiales y la inspección.}

\section{Una selección condicionada por el servicio}
Un material se selecciona para una pieza y unas condiciones de funcionamiento. La temperatura, el estado de tensiones, el ambiente, la geometría, el proceso de fabricación y la posibilidad de inspección modifican la comparación entre alternativas. Una propiedad aislada no basta para ordenar todas las soluciones.

En un motor, las exigencias cambian entre el fan, el compresor, la cámara de combustión y la turbina. El aumento de temperatura puede alterar las propiedades mecánicas y activar mecanismos que resultan poco relevantes a temperatura ambiente. La exposición ambiental también puede modificar la superficie y favorecer la iniciación del daño.

\caja{Propiedad, mecanismo y criterio}{La \textbf{propiedad} caracteriza una respuesta bajo condiciones definidas. El \textbf{mecanismo} explica cómo se produce la deformación o el daño. El \textbf{criterio} establece cómo se compara esa respuesta con una exigencia del problema. Deben declararse las condiciones de medida y los límites de uso.}

\section{Del daño a la decisión de inspección}
La inspección debe responder a una pregunta concreta: qué defecto se busca, en qué material, con qué orientación y en qué zona accesible. Las capacidades de una técnica dependen de esas condiciones. Una indicación requiere interpretación; su ausencia tampoco demuestra por sí sola la ausencia de cualquier daño.

\begin{figure}[htbp]
\centering
\begin{tikzpicture}[x=1cm,y=1cm,every node/.style={font=\small,align=center}]
\node[draw=studiumBlue,fill=studiumLight,text width=3.4cm,minimum height=1.25cm] (a) at (0,0) {Condiciones de servicio\\temperatura, carga, ambiente};
\node[draw=studiumBlue,fill=studiumLight,text width=3.4cm,minimum height=1.25cm] (b) at (5,0) {Respuesta del material\\deformación y degradación};
\node[draw=studiumBlue,fill=studiumLight,text width=3.4cm,minimum height=1.25cm] (c) at (0,-2.1) {Decisión de seguimiento\\datos y documentación};
\node[draw=studiumBlue,fill=studiumLight,text width=3.4cm,minimum height=1.25cm] (d) at (5,-2.1) {Evidencia de inspección\\método y limitaciones};
\draw[->,thick,studiumBlue] (a)--(b);\draw[->,thick,studiumBlue] (b)--(d);
\draw[->,thick,studiumBlue] (d)--(c);\draw[->,thick,studiumBlue] (c)--(a);
\end{tikzpicture}
\caption{Relación conceptual entre servicio, respuesta del material y seguimiento. Esquema original; no representa un procedimiento aprobado de mantenimiento.}
\label{fig:servicio}
\end{figure}

La figura~\ref{fig:servicio} relaciona la parte de materiales con la de mantenimiento. La cadena sirve para organizar el estudio: el mecanismo esperado orienta la búsqueda de evidencia y la evidencia permite revisar las hipótesis iniciales.

\clearpage
\section{Preguntas que ordenan el análisis}
El cuadro~\ref{tab:preguntas} reúne preguntas que deben acompañar al modelo. La composición conserva el tratamiento de cuadros de la colección y admite texto explicativo sin reducir el tamaño de letra.

\begin{table}[htbp]
\centering\small
\caption{Preguntas de partida para un análisis académico.}\label{tab:preguntas}
\begin{tabular}{p{.23\linewidth}p{.31\linewidth}p{.35\linewidth}}
\toprule
Aspecto & Pregunta & Información necesaria\\
\midrule
Carga & ¿Cómo varía con el tiempo? & Tensiones máximas y mínimas; secuencia de ciclos.\\[4pt]
Temperatura & ¿Es uniforme y constante? & Temperaturas, gradientes y tiempo de exposición.\\[4pt]
Ambiente & ¿Qué agentes actúan sobre la superficie? & Humedad, especies químicas y condiciones de contacto.\\[4pt]
Geometría & ¿Dónde se concentra el esfuerzo? & Radios, taladros, entallas y orientación de defectos.\\[4pt]
Inspección & ¿Qué defecto puede detectar la técnica? & Acceso, material, orientación, resolución y procedimiento.\\
\bottomrule
\end{tabular}
\end{table}

\subsection{Interpretación de una comprobación}
Una comprobación algebraica demuestra una identidad dentro de un modelo; una sustitución numérica reproduce una operación; un balance dimensional contrasta la coherencia de las unidades. Ninguna de estas comprobaciones demuestra, por sí sola, que el modelo sea adecuado para una pieza real.

Por ejemplo, una ley elástica lineal puede ser apropiada para una probeta sometida a pequeñas deformaciones. Su aplicación a una zona plastificada requiere revisar la hipótesis constitutiva, aunque las unidades de la ecuación sean correctas y el cálculo numérico no tenga errores.

\consejo{Pauta de estudio}{Antes de calcular, escriba el modelo y sus hipótesis. Después del cálculo, explique qué representa el resultado, cómo podría comprobarse y qué información faltaría para extenderlo al servicio real.}

\resumen{Las condiciones de servicio determinan qué propiedades y mecanismos deben analizarse. La selección del método de inspección requiere conocer el defecto buscado y las limitaciones de detección. El nivel de evidencia de cada comprobación debe declararse.}

\chapter{Tensión, deformación y rigidez}
\label{ch:elasticidad}
\objetivos{Distinguir fuerza, tensión y deformación; utilizar un modelo uniaxial lineal y comprobar las dimensiones de sus resultados.}

\section{Magnitudes y modelo uniaxial}
Para una barra recta de sección uniforme sometida a fuerza axial centrada, suficientemente lejos de la aplicación de la carga, la tensión nominal se expresa como
\begin{equation}\sigma=\frac{F}{A}.\label{eq:tension}\end{equation}
La deformación longitudinal media se define mediante la elongación y la longitud inicial:
\begin{equation}\varepsilon=\frac{\Delta L}{L_0}.\label{eq:deformacion}\end{equation}
En el régimen elástico lineal uniaxial se adopta
\begin{equation}\sigma=E\varepsilon,\qquad \Delta L=\frac{FL_0}{AE}.\label{eq:hooke}\end{equation}
El módulo $E$ relaciona tensión y deformación en este modelo. No equivale a una resistencia de rotura. Dos materiales pueden tener módulos próximos y resistencias diferentes.

\caja{Dominio de aplicación}{Las ecuaciones~\eqref{eq:tension}--\eqref{eq:hooke} se utilizan aquí con pequeñas deformaciones, sección uniforme, carga centrada y respuesta elástica lineal. No incluyen concentraciones locales de tensión, plasticidad, fluencia ni efectos térmicos.}

\section{Ejemplo desarrollado}
\begin{studiumworked}{E01. Elongación de una probeta idealizada}
\textbf{Enunciado y datos.} Una barra tiene $L_0=200\,\mathrm{mm}$, $A=50\,\mathrm{mm^2}$ y $E=70\,\mathrm{GPa}$. Se aplica $F=5{,}0\,\mathrm{kN}$. Calcule la tensión nominal y la elongación con el modelo anterior.

\textbf{Hipótesis.} Se supone que la respuesta permanece dentro del régimen elástico lineal; el enunciado no aporta un límite elástico con el que verificar esa condición material.

\textbf{Desarrollo y sustitución.} Se convierten primero las unidades a SI:
\begin{align*}
 A&=50\times10^{-6}\,\mathrm{m^2},&L_0&=0{,}200\,\mathrm{m},\\
 \sigma&=\frac{5000}{50\times10^{-6}}=1{,}0\times10^8\,\mathrm{Pa},\\
 \Delta L&=\frac{5000(0{,}200)}{(50\times10^{-6})(70\times10^9)}
 =2{,}857\times10^{-4}\,\mathrm{m}.
\end{align*}
\textbf{Resultado.} $\sigma=100\,\mathrm{MPa}$ y $\Delta L\simeq0{,}286\,\mathrm{mm}$.
\end{studiumworked}

\clearpage
\section{Comprobación e interpretación}
La deformación se obtiene por dos caminos independientes de sustitución:
\begin{equation}
\varepsilon=\frac{\sigma}{E}=\frac{100\times10^6}{70\times10^9}
=1{,}429\times10^{-3}=\frac{0{,}2857}{200}.
\label{eq:check}
\end{equation}
La igualdad de los valores comprueba la coherencia aritmética entre las ecuaciones elegidas. No verifica la respuesta constitutiva de una aleación concreta ni acredita una pieza para su uso.

\begin{tcolorbox}[studiumcommon,colback=studiumCheckBack,colframe=studiumCheckFrame,colbacktitle=studiumCheckFrame,coltitle=white,title={Comprobación dimensional}]
Para la tensión, $[F/A]=\mathrm{N/m^2}=\mathrm{Pa}$. Para la elongación,
\[
\left[\frac{FL_0}{AE}\right]
=\frac{\mathrm{N}\,\mathrm{m}}{\mathrm{m^2}\,\mathrm{N/m^2}}=\mathrm{m}.
\]
El resultado de la deformación es adimensional. La coherencia dimensional es necesaria, pero no suficiente, para validar el modelo físico.
\end{tcolorbox}

\subsection{Qué cambia al modificar un dato}
Dentro de este modelo, duplicar $F$ duplica tanto la tensión nominal como la elongación. Duplicar $A$ reduce ambas a la mitad. En cambio, duplicar $L_0$ duplica la elongación sin modificar la tensión nominal. Esta diferencia ayuda a separar una magnitud local idealizada de una respuesta geométrica global.

\subsection{Información que falta}
La temperatura de ensayo, el límite elástico, el estado metalúrgico y las posibles entallas no se han especificado. Por ello, el ejemplo no permite estimar vida a fatiga, margen frente a fluencia ni aceptabilidad de un daño. Para cualquiera de esos objetivos es necesario introducir datos y modelos adicionales.

\consejo{Comprobación antes de entregar}{Convierta el área con el cuadrado de la unidad de longitud. Si $1\,\mathrm{mm}=10^{-3}\,\mathrm{m}$, entonces $1\,\mathrm{mm^2}=10^{-6}\,\mathrm{m^2}$.}

\resumen{La fuerza se mide en newtons; la tensión, en pascales; la deformación es adimensional. El módulo elástico relaciona tensión y deformación bajo hipótesis concretas. Una solución numérica coherente debe acompañarse de sus límites de aplicación.}

\appendix
\chapter{Hoja de fórmulas de la muestra}
\label{ap:formulas}
\begin{table}[h]\centering
\caption{Relaciones utilizadas en el capítulo~\ref{ch:elasticidad}.}
\begin{tabular}{lll}\toprule
Magnitud & Relación & Unidad SI\\\midrule
Tensión nominal & $\sigma=F/A$ & $\mathrm{Pa}$\\[7pt]
Deformación media & $\varepsilon=\Delta L/L_0$ & $1$\\[7pt]
Ley uniaxial lineal & $\sigma=E\varepsilon$ & $\mathrm{Pa}$\\[7pt]
Elongación & $\Delta L=FL_0/(AE)$ & $\mathrm{m}$\\\bottomrule
\end{tabular}\end{table}
\caja{Alcance del formulario}{Estas relaciones pertenecen a una idealización uniaxial elástica. Deben leerse junto con las hipótesis del capítulo~\ref{ch:elasticidad}; no constituyen un formulario completo de la asignatura.}
